\section{Discussion}
\label{sec:discussion}

\paragraph{What the criterion decided}
Of the three clauses fixed before training, the first two hold and the third does not. The factored model beats the matched direct regressor on both held-out subject pairs and on both seeds, by 1.6 to 1.8\,mm with two views and by 4.4\,mm with one; the gain is concentrated where the formulation predicts, in far-field, occluded and out-of-view joints, and the far-field stratum carries the largest gain in every single-dataset comparison. It does not transfer to a second dataset, in any of four settings, and it disappears when a small amount of that dataset is added to training. The hypothesis that a factored representation helps is therefore not rejected, but its scope is narrower than the one we set out to test: factoring helps a regressor trained on one distribution of objects and poses, and the help is largest where the direct mapping is least constrained.

\paragraph{The gain is in the representation, not in the fusion}
The analysis of \cref{sec:inside} changes how the result should be read. We designed \method{} as two estimators combined by their precisions, expecting the geometric head to take over where the direct head is unsure. What happened is different. The geometric head, trained on the same tokens, is 11\,mm worse than the direct head; its predicted scale saturated on fold A; the fusion weight on it is a few percent; and the fused output is the direct head to within 0.2\,mm on every fold and seed. Yet the direct head inside the factored model is 1.4 to 4.5\,mm better than the same head trained without the geometric branch. The mechanism is auxiliary supervision through structure: forcing the shared decoder to produce joints and control points from which a differentiable operator can rebuild the field gives the direct head a better representation to regress from. This is a more modest claim than the one the architecture suggests, and a more useful one, because it identifies the component that can be removed at inference time. A practitioner who wants the gain without the 1.8-fold inference cost of the soft nearest-vertex operator can train with the geometric branch and deploy the direct head alone.

\paragraph{Why the fusion did not add more}
The per-joint oracle that picks the better head reaches 40.8\,mm on fold A, 3.7\,mm below what the fusion achieved; the geometric head is the better one on 28\% of joints. The complementarity exists and was not exploited, for two connected reasons. The geometric head's scale collapsed to a constant on fold A, so the fusion weight carried information about the direct head only; and where the direct head is uncertain, the geometric head is usually wrong as well (\cref{fig:uncertainty}c), so the weight signals difficulty rather than a preference. On fold B, where both scales trained normally and the geometric weight has a median of 0.13, the fusion still added only 0.2\,mm. Realising the oracle gap would require a geometric head that is right on a different set of joints from the direct head, not merely a better-calibrated one, which points at the object-pose branch: with the true object pose the geometric head is within 2\,mm of the direct head.

\paragraph{What limits the geometric route}
Two numbers bound it. In domain, replacing the predicted object pose by the truth lowers the geometric head's error by 8.7\,mm and replacing the predicted joints lowers it by 3.8\,mm; the predicted joints are already within the noise of the dataset's own hand annotation. Out of domain the ordering reverses, and the joint branch, not the object branch, is what fails first. A better object-pose branch is where the geometric formulation has headroom on SHOW3D; a joint branch that generalises across capture systems is what it would need to transfer. Neither is specific to interaction fields; both are the standard problems of hand-object pose estimation~\cite{fan2024benchmarks,bansal2026hopformer}, which the factored formulation inherits along with their solutions.

\paragraph{Relation to the challenge systems}
Our stereo gain of 3.4\,mm is eighteen times the 0.19\,mm GOLF measured at a LoRA-adapted ViT-H+, and our ray-encoding gain of 1.3\,mm is twice its 0.6\,mm~\cite{zou2026golf}. Our error components show that even at our representation the residual is mostly lateral rather than along the ray. Both observations say that the value of geometric structure shrinks as the backbone improves, because a stronger representation already encodes it. We expect the same of the factored gain: at the scale of the winning entries it may be smaller than the 1.7\,mm reported here, and an unadapted ViT-B is the setting in which it is easiest to measure. The result is a controlled measurement of one design decision, not a leaderboard entry, and the absolute numbers should be read accordingly.

\paragraph{Where the remaining error is}
Independently of the method, the analysis makes three statements about the benchmark that any future system will face. First, the mean is decided in the tail: 8\% of joints carry 43\% of the error, and the worst of them are hands that are not interacting with the object at all, for which the label is a vector of up to 1.3\,m. Second, one object handled by one subject accounts for a sixth of the fold-A error, so subject-level splits on this release test object generalisation as much as subject generalisation. Third, at contact both models are barely better than predicting the zero vector, and the object annotation noise is a floor there. A benchmark that separated interacting from non-interacting hands, or reported the field only for hands within some distance of the object, would measure the quantity its name promises more directly; the stratified numbers we release allow that separation post hoc.
